\documentclass[12pt]{article}
\usepackage[utf8]{inputenc}
\usepackage[T1]{fontenc}
\usepackage{hyperref}
\title{Pathiya Language Resources for K-12 Curriculum Development}
\author{Compiled by Claude Sonnet 5.5 (claude-sonnet-5-5) for Andr\'as Kornai}
\date{}
\begin{document}
\maketitle

\section{Introduction}
Pathiya (ISO 639-3: pty; Glottocode path1235) is a small Southern Dravidian language of India. A web search found only encyclopedic coverage: it is described as close to Kalanadi (about 88\% lexical similarity) and Kunduvadi, with roughly 1,000 speakers reported in 2004. Searches turned up no primers, textbooks, dictionaries, audio or school programmes for Pathiya. Note that Paniya (pcg), a different language of Wayanad, appeared in the same results and must not be confused with Pathiya.

\section{Primary Educational Resources}
None found. No primer, textbook or school programme for Pathiya was located in search results.

\section{Linguistic and Reference Materials}
\begin{itemize}
\item Pathiya language (encyclopedia overview, classification and related languages): \url{https://en.wikipedia.org/wiki/Pathiya_language}
\item Kalanadi language (closest reported relative): \url{https://en.wikipedia.org/wiki/Kalanadi_language}
\item Kunduvadi language (related Southern Dravidian language): \url{https://en.wikipedia.org/wiki/Kunduvadi_language}
\end{itemize}

\section{Religious and Community Resources}
None found.

\section{Digital and Audio Resources}
None found.

\section{Curriculum Development Context}
Little is documented about Pathiya. Any curriculum work would first require primary documentation with the community (word lists, texts, recordings). The only materials located are general reference pages, which are not teaching resources. Related languages such as Kalanadi and Kunduvadi are likewise sparsely documented.
\end{document}